\section{Institutional Background} \label{sec_background}

The Los Angeles planning and approvals process for urban development is a multi-layered process that requires the input of multiple agencies. If a development is ``by-right'', meaning that it conforms to all zoning regulations within its applicable zone, then development can proceed with a permit from the Department of Building and Safety. If, however, the planned development requires an exception to the zoning code, or has a planned purpose not normally allowed by the zoning code, then approval needs to be obtained from the Planning Department and potentially from the City Planning Commission (CPC).

The LA CPC is a nine member administrative body with members appointed by the Mayor of LA and confirmed by the City Council. The CPC is responsible for reviewing citywide zoning changes, handling approvals and conditional use permits for large developments, and handling appeals of decisions made by lower bodies. The LA CPC is the highest planning body within the City's Planning Department, and sits below only the City Council in the City's overall planning hierarchy. CPC appellate decisions are often final and not appealable, and when they are appealable they may only be appealed to the City Council. 

The CPC meets approximately every two weeks. The agenda for each meeting is set by the Commission itself and made available to the public at least seven days before each meeting. After each meeting, the meeting minutes are published online. The minutes record the order of discussion, the members present, the motions made on each agenda item, and the votes made on each motion. Motions are passed by majority rule, subject to a quorum of five members present. In addition, members of the public are allowed to submit comments on any agenda item either in writing or in person during the meeting. 
